% answers.tex —— 2012 年·广东卷·文科 · 参考答案与详细解析片段
% 由 tools/generate_projects.py 依来源总账自 2012/2012_数学真题_广东文_试题.md + 2012/2012_数学真题_广东文_解析正文.md + 2012/2012_数学真题_广东文_答案要点.md 生成，勿手改（改源 MD 后重跑）


\begin{answer}
\gksolhead{1}{}

\textbf{题目}：

设 $\mathrm i$ 为虚数单位，则复数 $\dfrac{3+4\mathrm i}{\mathrm i}=$\mbox{（　　）}

（A）$-4-3\mathrm i$；（B）$-4+3\mathrm i$；（C）$4+3\mathrm i$；（D）$4-3\mathrm i$。

\Needspace{4\baselineskip}
\textbf{解析}

\Needspace{5\baselineskip}
由 $\dfrac1{\mathrm i}=\dfrac{\mathrm i}{\mathrm i^2}=-\mathrm i$，把除法化为乘法：
\nopagebreak[4]
\[\frac{3+4\mathrm i}{\mathrm i}=(3+4\mathrm i)\cdot(-\mathrm i)=-3\mathrm i-4\mathrm i^2.\]
\Needspace{5\baselineskip}
代入 $\mathrm i^2=-1$，得
\nopagebreak[4]
\[-3\mathrm i+4=4-3\mathrm i.\]
实部为 $4$、虚部为 $-3$。

选 D。

\gkanswer{答案：D。$\dfrac1{\mathrm i}=-\mathrm i$，故 $\dfrac{3+4\mathrm i}{\mathrm i}=(3+4\mathrm i)(-\mathrm i)=4-3\mathrm i$。}

\end{answer}
\begin{answer}
\gksolhead{2}{}

\textbf{题目}：

设集合 $U=\{1,2,3,4,5,6\}$，$M=\{1,3,5\}$，则 $\complement_U M=$\mbox{（　　）}

（A）$\{2,4,6\}$；（B）$\{1,3,5\}$；（C）$\{1,2,4\}$；（D）$U$。

\Needspace{4\baselineskip}
\textbf{解析}

\Needspace{5\baselineskip}
$\complement_U M$ 由 $U$ 中不属于 $M$ 的元素组成。逐个核对 $U=\{1,2,3,4,5,6\}$：
\nopagebreak[4]
\[1,3,5\in M;\qquad 2,4,6\notin M.\]
\Needspace{5\baselineskip}
故
\nopagebreak[4]
\[\complement_U M=\{2,4,6\}.\]

选 A。

\gkanswer{答案：A。$U$ 中去掉 $M=\{1,3,5\}$ 的元素，剩 $\{2,4,6\}$。}

\end{answer}
\begin{answer}
\gksolhead{3}{}

\textbf{题目}：

若向量 $\overrightarrow{AB}=(1,2)$，$\overrightarrow{BC}=(3,4)$，则 $\overrightarrow{AC}=$\mbox{（　　）}

（A）$(4,6)$；（B）$(-4,-6)$；（C）$(-2,-2)$；（D）$(2,2)$。

\Needspace{4\baselineskip}
\textbf{解析}

\Needspace{5\baselineskip}
由向量加法的三角形法则，
\nopagebreak[4]
\[\overrightarrow{AC}=\overrightarrow{AB}+\overrightarrow{BC}.\]
\Needspace{5\baselineskip}
按坐标相加：
\nopagebreak[4]
\[\overrightarrow{AC}=(1,2)+(3,4)=(1+3,\,2+4)=(4,6).\]

选 A。

\gkanswer{答案：A。$\overrightarrow{AC}=\overrightarrow{AB}+\overrightarrow{BC}=(1+3,2+4)=(4,6)$。}

\end{answer}
\begin{answer}
\gksolhead{4}{}

\textbf{题目}：

下列函数为偶函数的是\mbox{（　　）}

（A）$y=\sin x$；（B）$y=x^3$；（C）$y=\mathrm e^x$；（D）$y=\ln\sqrt{x^2+1}$。

\Needspace{4\baselineskip}
\textbf{解析}

偶函数要求定义域关于原点对称，且对定义域内每个 $x$ 都有 $f(-x)=f(x)$。逐项检验：

\begin{itemize}
\item $y=\sin x$：$\sin(-x)=-\sin x$，是奇函数；
\item $y=x^3$：$(-x)^3=-x^3$，是奇函数；
\item $y=\mathrm e^x$：$\mathrm e^{-x}=\dfrac1{\mathrm e^x}\neq\mathrm e^x$，非奇非偶；
\item $y=\ln\sqrt{x^2+1}$：定义域为 $\mathbf R$，且
\end{itemize}
\nopagebreak[4]
\[\ln\sqrt{(-x)^2+1}=\ln\sqrt{x^2+1},\]
满足 $f(-x)=f(x)$，是偶函数。

选 D。

\gkanswer{答案：D。$\sin x$ 与 $x^3$ 为奇函数，$\mathrm e^x$ 非奇非偶；只有 $\ln\sqrt{x^2+1}$ 满足 $f(-x)=f(x)$。}

\end{answer}
\begin{answer}
\gksolhead{5}{}

\textbf{题目}：

\Needspace{5\baselineskip}
已知变量 $x$，$y$ 满足约束条件
\nopagebreak[4]
\[\begin{cases}x+y\leqslant1,\\ x+1\geqslant0,\\ x-y\leqslant1,\end{cases}\]
则 $z=x+2y$ 的最小值为\mbox{（　　）}

（A）$3$；（B）$1$；（C）$-5$；（D）$-6$。

\Needspace{4\baselineskip}
\textbf{解析}

\Needspace{5\baselineskip}
三条边界直线为
\nopagebreak[4]
\[l_1:x+y=1,\qquad l_2:x=-1,\qquad l_3:y=x-1.\]
\Needspace{5\baselineskip}
两两求交，得可行域的三个顶点
\nopagebreak[4]
\[(-1,2),\qquad (-1,-2),\qquad (1,0).\]
\Needspace{5\baselineskip}
逐点计算 $z=x+2y$：
\nopagebreak[4]
\[\begin{aligned}
(-1,2)&:\quad z=-1+2\times2=3,\\
(-1,-2)&:\quad z=-1+2\times(-2)=-5,\\
(1,0)&:\quad z=1+2\times0=1.
\end{aligned}\]
线性目标函数在可行域的顶点处取得最值，故 $z_{\min}=-5$。

选 C。

\gkanswer{答案：C。可行域三个顶点为 $(-1,2)$，$(-1,-2)$，$(1,0)$，对应 $z$ 为 $3$，$-5$，$1$，最小值 $-5$。}

\end{answer}
\begin{answer}
\gksolhead{6}{}

\textbf{题目}：

在 $\triangle ABC$ 中，若 $\angle A=60^\circ$，$\angle B=45^\circ$，$BC=3\sqrt2$，则 $AC=$\mbox{（　　）}

（A）$4\sqrt3$；（B）$2\sqrt3$；（C）$\sqrt3$；（D）$\dfrac{\sqrt3}{2}$。

\Needspace{4\baselineskip}
\textbf{解析}

\Needspace{5\baselineskip}
在 $\triangle ABC$ 中，$BC$ 是 $\angle A$ 的对边，$AC$ 是 $\angle B$ 的对边。由正弦定理
\nopagebreak[4]
\[\frac{AC}{\sin B}=\frac{BC}{\sin A}.\]
\Needspace{5\baselineskip}
代入 $B=45^\circ$，$A=60^\circ$，$BC=3\sqrt2$：
\nopagebreak[4]
\[AC=\frac{BC\sin B}{\sin A}=\frac{3\sqrt2\cdot\frac{\sqrt2}{2}}{\frac{\sqrt3}{2}}=\frac{3}{\frac{\sqrt3}{2}}=\frac{6}{\sqrt3}=2\sqrt3.\]

选 B。

\gkanswer{答案：B。正弦定理 $\dfrac{AC}{\sin45^\circ}=\dfrac{BC}{\sin60^\circ}$，得 $AC=\dfrac{3\sqrt2\cdot\frac{\sqrt2}2}{\frac{\sqrt3}2}=2\sqrt3$。}

\end{answer}
\begin{answer}
\gksolhead{7}{}

\textbf{题目}：

某几何体的三视图如图所示，则它的体积为\mbox{（　　）}
\nopagebreak[4]
\begin{center}% ============================================================
% f2012_guangdong_arts_q07_views.tex
% 2012 广东卷（文科）第 7 题 三视图（TikZ 矢量重绘）
% 正视图、侧视图：上半为半径 3 的半圆，下半为底 6、腰 5（高 4）的倒三角
% 俯视图：直径 6 的圆；还原为半球（r=3）与圆锥（r=3, h=4）的组合体
% ============================================================
\begin{tikzpicture}[
  x=1cm, y=1cm, line width=0.6pt, font=\footnotesize, >=Stealth
]

  % ---------------- 正视图 ----------------
  \begin{scope}
    % 半圆（上）
    \draw (-1.5,0) arc[start angle=180, end angle=0, radius=1.5];
    \draw (-1.5,0) -- (1.5,0);
    % 三角形（下，顶点朝下）
    \draw (-1.5,0) -- (0,-2) -- (1.5,0);
    % 上方尺寸 6
    \draw[<->] (-1.5,1.85) -- (1.5,1.85);
    \draw (-1.5,1.78) -- (-1.5,1.92);
    \draw (1.5,1.78) -- (1.5,1.92);
    \node[above] at (0,1.90) {$6$};
    % 右侧尺寸 3
    \draw[<->] (1.78,0) -- (1.78,1.5);
    \draw (1.72,0) -- (1.86,0);
    \draw (1.72,1.5) -- (1.86,1.5);
    \node[right] at (1.80,0.75) {$3$};
    % 两腰尺寸 5
    \node[rotate=-53.13] at (-1.12,-1.24) {$5$};
    \node[rotate=53.13] at (1.12,-1.24) {$5$};
    % 视图名
    \node at (0,-2.55) {正视图};
  \end{scope}

  % ---------------- 侧视图 ----------------
  \begin{scope}[xshift=4.9cm]
    \draw (-1.5,0) arc[start angle=180, end angle=0, radius=1.5];
    \draw (-1.5,0) -- (1.5,0);
    \draw (-1.5,0) -- (0,-2) -- (1.5,0);
    \draw[<->] (-1.5,1.85) -- (1.5,1.85);
    \draw (-1.5,1.78) -- (-1.5,1.92);
    \draw (1.5,1.78) -- (1.5,1.92);
    \node[above] at (0,1.90) {$6$};
    \draw[<->] (1.78,0) -- (1.78,1.5);
    \draw (1.72,0) -- (1.86,0);
    \draw (1.72,1.5) -- (1.86,1.5);
    \node[right] at (1.80,0.75) {$3$};
    \node[rotate=-53.13] at (-1.12,-1.24) {$5$};
    \node[rotate=53.13] at (1.12,-1.24) {$5$};
    \node at (0,-2.55) {侧视图};
  \end{scope}

  % ---------------- 俯视图 ----------------
  \begin{scope}[yshift=-5.8cm]
    \draw (0,0) circle (1.5);
    \node at (0,-1.95) {俯视图};
  \end{scope}

\end{tikzpicture}
\end{center}
\nopagebreak[4]
（A）$72\pi$；（B）$48\pi$；（C）$30\pi$；（D）$24\pi$。

\Needspace{4\baselineskip}
\textbf{解析}

正视图由上半部的半圆和下半部的等腰三角形拼成，侧视图与正视图全同，俯视图是一个圆（直径 $6$），说明该几何体由两部分组成：

\begin{itemize}
\item 上半部是半径为 $3$ 的半球，高为 $3$；
\item 下半部是底面半径 $3$、母线 $5$ 的圆锥，其高
\end{itemize}
\nopagebreak[4]
\[h=\sqrt{5^2-3^2}=\sqrt{25-9}=4.\]

\Needspace{5\baselineskip}
于是
\nopagebreak[4]
\[\begin{aligned}
V_{\text{\text{半球}}}&=\frac23\pi\cdot3^3=18\pi,\\
V_{\text{\text{圆锥}}}&=\frac13\pi\cdot3^2\cdot4=12\pi.
\end{aligned}\]
\Needspace{5\baselineskip}
合计
\nopagebreak[4]
\[V=18\pi+12\pi=30\pi.\]

选 C。

\gkanswer{答案：C。几何体为半球与圆锥的组合，半径均为 $3$，锥高 $\sqrt{5^2-3^2}=4$，故 $V=\dfrac23\pi\cdot27+\dfrac13\pi\cdot9\cdot4=30\pi$。}

\end{answer}
\begin{answer}
\gksolhead{8}{}

\textbf{题目}：

在平面直角坐标系 $xOy$ 中，直线 $3x+4y-5=0$ 与圆 $x^2+y^2=4$ 相交于 $A$，$B$ 两点，则弦 $AB$ 的长等于\mbox{（　　）}

（A）$3\sqrt3$；（B）$2\sqrt3$；（C）$\sqrt3$；（D）$1$。

\Needspace{4\baselineskip}
\textbf{解析}

\Needspace{5\baselineskip}
圆 $x^2+y^2=4$ 的圆心为 $O(0,0)$，半径 $r=2$。圆心到直线 $3x+4y-5=0$ 的距离
\nopagebreak[4]
\[d=\frac{|3\times0+4\times0-5|}{\sqrt{3^2+4^2}}=\frac{5}{5}=1.\]
\Needspace{5\baselineskip}
弦长公式给出
\nopagebreak[4]
\[|AB|=2\sqrt{r^2-d^2}=2\sqrt{4-1}=2\sqrt3.\]

选 B。

\gkanswer{答案：B。圆心到直线距离 $d=\dfrac{|{-5}|}{5}=1$，故 $|AB|=2\sqrt{4-1^2}=2\sqrt3$。}

\end{answer}
\begin{answer}
\gksolhead{9}{}

\textbf{题目}：

执行如图所示的程序框图，若输入 $n$ 的值为 $6$，则输出 $s$ 的值为\mbox{（　　）}
\nopagebreak[4]
\begin{center}% ============================================================
% f2012_guangdong_arts_q09_flowchart.tex
% 2012 广东卷（文科）第 9 题 程序框图（TikZ 矢量重绘）
% 开始 → 输入 n → i=1, s=1 → 判断 i<n
%   是：s=s×i → i=i+2 → 回到判断
%   否：输出 s → 结束
% ============================================================
\begin{tikzpicture}[
  x=1cm, y=1cm, line width=0.6pt, font=\footnotesize, >=Stealth,
  every node/.style={inner sep=3pt},
  proc/.style={rectangle, draw, minimum width=2.5cm, minimum height=0.62cm},
  io/.style={trapezium, draw, trapezium left angle=70, trapezium right angle=110,
             minimum width=2.7cm, minimum height=0.66cm, trapezium stretches=true},
  dec/.style={diamond, draw, aspect=2.4, minimum width=2.5cm, minimum height=0.9cm},
  term/.style={rectangle, rounded corners=8pt, draw, minimum width=1.9cm, minimum height=0.58cm}
]

  \node[term] (start) at (0,0) {开始};
  \node[io]   (input) at (0,-1.30) {输入 $n$};
  \node[proc] (init)  at (0,-2.60) {$i=1$，$s=1$};
  \node[dec]  (judge) at (0,-4.15) {$i<n$};

  \node[proc] (mul)   at (0,-5.65) {$s=s\times i$};
  \node[proc] (inc)   at (0,-6.95) {$i=i+2$};

  \node[io]   (out)   at (3.50,-5.65) {输出 $s$};
  \node[term] (end)   at (3.50,-6.95) {结束};

  \draw[->] (start) -- (input);
  \draw[->] (input) -- (init);
  \draw[->] (init) -- (judge);
  \draw[->] (judge) -- node[right] {是} (mul);
  \draw[->] (mul) -- (inc);
  \draw[->] (inc.west) -- (-2.20,-6.95) -- (-2.20,-4.15) -- (judge.west);
  \draw[->] (judge.east) -- (3.50,-4.15) -- (out.north);
  \node at (2.30,-3.95) {否};
  \draw[->] (out) -- (end);

\end{tikzpicture}
\end{center}
\nopagebreak[4]
（A）$105$；（B）$16$；（C）$15$；（D）$1$。

\Needspace{4\baselineskip}
\textbf{解析}

\Needspace{5\baselineskip}
按框图逐步执行，初值 $i=1$，$s=1$，输入 $n=6$：

\nopagebreak[4]
\[\begin{aligned}
i=1<6&:\quad s=1\times1=1,\quad i=1+2=3;\\
i=3<6&:\quad s=1\times3=3,\quad i=3+2=5;\\
i=5<6&:\quad s=3\times5=15,\quad i=5+2=7;\\
i=7<6&\text{ \text{不成立，输出} } s=15.
\end{aligned}\]

选 C。

\gkanswer{答案：C。依次为 $s=1,i=3$；$s=3,i=5$；$s=15,i=7$，此时 $7<6$ 不成立，输出 $15$。}

\end{answer}
\begin{answer}
\gksolhead{10}{}

\textbf{题目}：

对任意两个非零的平面向量 $\boldsymbol\alpha$ 和 $\boldsymbol\beta$，定义 $\boldsymbol\alpha\circ\boldsymbol\beta=\dfrac{\boldsymbol\alpha\cdot\boldsymbol\beta}{\boldsymbol\beta\cdot\boldsymbol\beta}$。若两个非零的平面向量 $\boldsymbol a$，$\boldsymbol b$ 满足 $\boldsymbol a$ 与 $\boldsymbol b$ 的夹角 $\theta\in\left(\dfrac\pi4,\dfrac\pi2\right)$，且 $\boldsymbol a\circ\boldsymbol b$ 和 $\boldsymbol b\circ\boldsymbol a$ 都在集合 $\left\{\dfrac n2\mid n\in\mathbf Z\right\}$ 中，则 $\boldsymbol a\circ\boldsymbol b=$\mbox{（　　）}

（A）$\dfrac52$；（B）$\dfrac32$；（C）$1$；（D）$\dfrac12$。

\Needspace{4\baselineskip}
\textbf{解析}

\Needspace{5\baselineskip}
记 $|\boldsymbol a|=a$，$|\boldsymbol b|=b$，两向量夹角为 $\theta$，则 $\boldsymbol a\cdot\boldsymbol b=ab\cos\theta$，$\boldsymbol b\cdot\boldsymbol b=b^2$，于是
\nopagebreak[4]
\[\boldsymbol a\circ\boldsymbol b=\frac{ab\cos\theta}{b^2}=\frac ab\cos\theta,\qquad
\boldsymbol b\circ\boldsymbol a=\frac{ab\cos\theta}{a^2}=\frac ba\cos\theta.\]
\Needspace{5\baselineskip}
两式相乘，$a$，$b$ 约去：
\nopagebreak[4]
\[(\boldsymbol a\circ\boldsymbol b)(\boldsymbol b\circ\boldsymbol a)=\cos^2\theta.\]

\Needspace{5\baselineskip}
因为 $\theta\in\left(\dfrac\pi4,\dfrac\pi2\right)$，所以 $0<\cos\theta<\dfrac{\sqrt2}{2}$，从而
\nopagebreak[4]
\[0<\cos^2\theta<\frac12.\]

\Needspace{5\baselineskip}
设 $\boldsymbol a\circ\boldsymbol b=\dfrac m2$，$\boldsymbol b\circ\boldsymbol a=\dfrac k2$（$m$，$k\in\mathbf Z$）。由 $\cos\theta>0$ 知两式均为正，故 $m\geqslant1$，$k\geqslant1$，且
\nopagebreak[4]
\[\frac{mk}{4}=\cos^2\theta<\frac12,\]
\Needspace{5\baselineskip}
即 $mk<2$。又 $mk\geqslant1$ 为整数，只能 $mk=1$，即 $m=k=1$。所以
\nopagebreak[4]
\[\boldsymbol a\circ\boldsymbol b=\frac12.\]

选 D。

\gkanswer{答案：D。两式之积 $(\boldsymbol a\circ\boldsymbol b)(\boldsymbol b\circ\boldsymbol a)=\cos^2\theta\in\left(0,\dfrac12\right)$，而两者均为 $\dfrac n2$（$n$ 为正整数），其积为 $\dfrac{mk}4$，只能 $mk=1$。}

\end{answer}
\begin{answer}
\gksolhead{11}{}

\textbf{题目}：

函数 $y=\dfrac{\sqrt{x+1}}{x}$ 的定义域为\gkblank{}。

\Needspace{4\baselineskip}
\textbf{解析}

\Needspace{5\baselineskip}
根号内非负、分母非零，故
\nopagebreak[4]
\[\begin{cases}x+1\geqslant0,\\ x\neq0,\end{cases}\]
\Needspace{5\baselineskip}
即 $x\geqslant-1$ 且 $x\neq0$。写成区间为
\nopagebreak[4]
\[[-1,0)\cup(0,+\infty).\]

\gkanswer{答案：$[-1,0)\cup(0,+\infty)$。需 $x+1\geqslant0$ 且 $x\neq0$。}

\end{answer}
\begin{answer}
\gksolhead{12}{}

\textbf{题目}：

等比数列 $\{a_n\}$ 满足 $a_2a_4=\dfrac12$，则 $a_1a_3^2a_5=$\gkblank{}。

\Needspace{4\baselineskip}
\textbf{解析}

\Needspace{5\baselineskip}
设公比为 $q$，则 $a_2=a_3q^{-1}$，$a_4=a_3q$，故
\nopagebreak[4]
\[a_2a_4=a_3^2.\]
\Needspace{5\baselineskip}
由已知 $a_2a_4=\dfrac12$ 得
\nopagebreak[4]
\[a_3^2=\frac12.\]
\Needspace{5\baselineskip}
同理 $a_1=a_3q^{-2}$，$a_5=a_3q^2$，故 $a_1a_5=a_3^2=\dfrac12$。于是
\nopagebreak[4]
\[a_1a_3^2a_5=(a_1a_5)\cdot a_3^2=\frac12\cdot\frac12=\frac14.\]

\gkanswer{答案：$\dfrac14$。$a_3^2=a_2a_4=\dfrac12$，$a_1a_5=a_3^2$，故原式 $=\left(\dfrac12\right)^2=\dfrac14$。}

\end{answer}
\begin{answer}
\gksolhead{13}{}

\textbf{题目}：

由正整数组成的一组数据 $x_1$，$x_2$，$x_3$，$x_4$，其平均数和中位数都是 $2$，且标准差等于 $1$，则这组数据为\gkblank{}。（从小到大排列）

\Needspace{4\baselineskip}
\textbf{解析}

设四个正整数从小到大为 $x_1\leqslant x_2\leqslant x_3\leqslant x_4$。

\Needspace{5\baselineskip}
平均数为 $2$ 给出
\nopagebreak[4]
\[x_1+x_2+x_3+x_4=8.\]
\Needspace{5\baselineskip}
中位数为 $2$ 给出
\nopagebreak[4]
\[\frac{x_2+x_3}{2}=2,\qquad x_2+x_3=4,\]
正整数解只有 $(x_2,x_3)=(2,2)$ 或 $(1,3)$。

\Needspace{5\baselineskip}
标准差为 $1$，即方差为 $1$：
\nopagebreak[4]
\[\frac14\sum_{i=1}^{4}(x_i-2)^2=1,\qquad \sum_{i=1}^{4}(x_i-2)^2=4.\]

\Needspace{5\baselineskip}
若 $(x_2,x_3)=(2,2)$，则 $(x_1-2)^2+(x_4-2)^2=4$，又 $x_1+x_4=4$，令 $x_4=4-x_1$ 得
\nopagebreak[4]
\[2(x_1-2)^2=4,\qquad (x_1-2)^2=2,\]
无整数解。

\Needspace{5\baselineskip}
若 $(x_2,x_3)=(1,3)$，则 $(x_1-2)^2+(x_4-2)^2=2$，同样代入 $x_4=4-x_1$：
\nopagebreak[4]
\[2(x_1-2)^2=2,\qquad (x_1-2)^2=1,\]
得 $x_1=1$ 或 $x_1=3$。由 $x_1\leqslant x_2=1$ 取 $x_1=1$，此时 $x_4=3$。

\Needspace{5\baselineskip}
故这组数据为
\nopagebreak[4]
\[1,1,3,3.\]
检验：平均数 $\dfrac{1+1+3+3}{4}=2$，中位数 $\dfrac{1+3}{2}=2$，方差 $\dfrac{4\times1}{4}=1$，标准差为 $1$，全部符合。

\gkanswer{答案：$1,1,3,3$。四数和为 $8$、中间两数和为 $4$，逐一检验方差为 $1$ 的组合，仅此一组。}

\end{answer}
\begin{answer}
\gksolhead{14}{}

\textbf{题目}：

\Needspace{5\baselineskip}
在平面直角坐标系 $xOy$ 中，曲线 $C_1$ 和 $C_2$ 的参数方程分别为
\nopagebreak[4]
\[\begin{cases}x=\sqrt5\cos\theta,\\ y=\sqrt5\sin\theta,\end{cases}\quad(\theta\text{ \text{为参数，}}0\leqslant\theta\leqslant\tfrac\pi2)\]
\Needspace{5\baselineskip}
和
\nopagebreak[4]
\[\begin{cases}x=1-\dfrac{\sqrt2}{2}t,\\ y=-\dfrac{\sqrt2}{2}t,\end{cases}\quad(t\text{ \text{为参数}}),\]
则曲线 $C_1$ 与 $C_2$ 的交点坐标为\gkblank{}。

\Needspace{4\baselineskip}
\textbf{解析}

\Needspace{5\baselineskip}
对 $C_1$，$\cos\theta=\dfrac{x}{\sqrt5}$，$\sin\theta=\dfrac{y}{\sqrt5}$，平方相加得
\nopagebreak[4]
\[x^2+y^2=5.\]
又 $0\leqslant\theta\leqslant\dfrac\pi2$，故 $x\geqslant0$，$y\geqslant0$，$C_1$ 是第一象限内的四分之一圆弧。

\Needspace{5\baselineskip}
对 $C_2$，由 $y=-\dfrac{\sqrt2}{2}t$ 得 $\sqrt2t=-2y$，即 $t=-\sqrt2y$。代入第一式：
\nopagebreak[4]
\[x=1-\frac{\sqrt2}{2}\cdot\left(-\sqrt2y\right)=1+y,\]
\Needspace{5\baselineskip}
即
\nopagebreak[4]
\[y=x-1.\]
故 $C_2$ 是直线 $y=x-1$。

\Needspace{5\baselineskip}
联立
\nopagebreak[4]
\[x^2+(x-1)^2=5,\]
\Needspace{5\baselineskip}
得
\nopagebreak[4]
\[2x^2-2x-4=0,\qquad x^2-x-2=0,\]
即 $x=2$ 或 $x=-1$。$x=-1$ 对应 $y=-2$，不满足 $x\geqslant0$，$y\geqslant0$，舍去；$x=2$ 对应 $y=1$，符合。

\Needspace{5\baselineskip}
故交点坐标为
\nopagebreak[4]
\[(2,1).\]
代回检验：$2^2+1^2=5$，且 $1=2-1$，两式同时成立。

\gkanswer{答案：$(2,1)$。$C_1$ 是圆 $x^2+y^2=5$ 在第一象限的四分之一圆弧，$C_2$ 是直线 $y=x-1$，联立得 $x=2$ 或 $x=-1$，舍去后者。}

\end{answer}
\begin{answer}
\gksolhead{15}{}

\textbf{题目}：

如图所示，直线 $PB$ 与圆 $O$ 相切于点 $B$，$D$ 是弦 $AC$ 上的点，$\angle PBA=\angle DBA$，若 $AD=m$，$AC=n$，则 $AB=$\gkblank{}。
\nopagebreak[4]
\begin{center}% ============================================================
% f2012_guangdong_arts_q15_circle.tex
% 2012 广东卷（文科）第 15 题（几何证明选讲）圆与弦（TikZ 矢量重绘）
% 圆 O 半径 1.8，B 在圆上，PB 切圆于 B，D 在弦 AC 上，∠PBA=∠DBA
% AD=m，AC=n，结论 AB=√(mn)；按原页拓扑示意，不按比例
% ============================================================
\begin{tikzpicture}[
  x=1cm, y=1cm, line width=0.6pt, font=\footnotesize, >=Stealth
]

  \coordinate (O) at (0,0);
  \coordinate (B) at (-1.691,-0.616);
  \coordinate (A) at (-0.466,1.739);
  \coordinate (C) at (1.691,-0.616);
  \coordinate (D) at (0.721,0.444);
  \coordinate (P) at (-2.444,1.452);

  % 圆
  \draw (O) circle (1.8);

  % 切线 PB（并向 B 外延长少许）
  \draw (P) -- (B) -- ++(0.342,-0.940) coordinate (Bt);
  % 弦 AC 与 D
  \draw (A) -- (C);
  % 三角形 ABC 的边
  \draw (A) -- (B);
  \draw (B) -- (C);
  % BD
  \draw (B) -- (D);

  % 圆心
  \fill (O) circle (0.035);

  % 标注
  \node[left] at (P) {$P$};
  \node[above left] at (A) {$A$};
  \node[below left] at (B) {$B$};
  \node[right] at (C) {$C$};
  \node[above right] at (D) {$D$};
  \node[below right] at (O) {$O$};

\end{tikzpicture}
\end{center}
\nopagebreak[4]
\Needspace{4\baselineskip}
\textbf{解析}

\Needspace{5\baselineskip}
由弦切角定理，弦切角等于它所夹弧上的圆周角，即
\nopagebreak[4]
\[\angle PBA=\angle BCA.\]
\Needspace{5\baselineskip}
已知 $\angle PBA=\angle DBA$，故
\nopagebreak[4]
\[\angle DBA=\angle BCA.\]
\Needspace{5\baselineskip}
又 $D$ 在弦 $AC$ 上，$\angle BAD$ 与 $\angle BAC$ 是同一个角（射线 $AD$ 与 $AC$ 重合），因此
\nopagebreak[4]
\[\angle BAD=\angle BAC.\]
\Needspace{5\baselineskip}
在 $\triangle ABD$ 与 $\triangle ACB$ 中，两组角分别相等，故
\nopagebreak[4]
\[\triangle ABD\sim\triangle ACB.\]
\Needspace{5\baselineskip}
对应边成比例（$AB$ 对应 $AC$，$AD$ 对应 $AB$）：
\nopagebreak[4]
\[\frac{AB}{AC}=\frac{AD}{AB},\]
\Needspace{5\baselineskip}
即
\nopagebreak[4]
\[AB^2=AD\cdot AC=mn,\]
\Needspace{5\baselineskip}
所以
\nopagebreak[4]
\[AB=\sqrt{mn}.\]

\gkanswer{答案：$\sqrt{mn}$。由弦切角定理及 $\angle PBA=\angle DBA$ 得 $\angle DBA=\angle BCA$，故 $\triangle ABD\backsim\triangle ACB$，$AB^2=AD\cdot AC=mn$。}

\end{answer}
\begin{answer}
\gksolhead{16}{}

\textbf{题目}：

已知函数 $f(x)=A\cos\left(\dfrac x4+\dfrac\pi6\right)$，$x\in\mathbf R$，且 $f\left(\dfrac\pi3\right)=\sqrt2$。

（1） 求 $A$ 的值；

（2） 设 $\alpha$，$\beta\in\left[0,\dfrac\pi2\right]$，$f\left(4\alpha+\dfrac{4\pi}{3}\right)=-\dfrac{30}{17}$，$f\left(4\beta-\dfrac{2\pi}{3}\right)=\dfrac85$，求 $\cos(\alpha+\beta)$ 的值。

\Needspace{4\baselineskip}
\textbf{解析}

\textbf{（1）}

\Needspace{5\baselineskip}
由 $f\left(\dfrac\pi3\right)=\sqrt2$，代入 $x=\dfrac\pi3$：
\nopagebreak[4]
\[\frac x4+\frac\pi6=\frac\pi{12}+\frac{\pi}{6}=\frac{\pi}{4},\]
\Needspace{5\baselineskip}
于是
\nopagebreak[4]
\[A\cos\frac\pi4=\sqrt2,\qquad A\cdot\frac{\sqrt2}{2}=\sqrt2,\]
\Needspace{5\baselineskip}
解得
\nopagebreak[4]
\[A=2.\]

\textbf{（2）}

\Needspace{5\baselineskip}
先化简两个条件。对第一个条件，自变量除以 $4$ 后加 $\dfrac\pi6$：
\nopagebreak[4]
\[\frac14\left(4\alpha+\frac{4\pi}{3}\right)+\frac\pi6=\alpha+\frac\pi3+\frac\pi6=\alpha+\frac\pi2,\]
\Needspace{5\baselineskip}
故
\nopagebreak[4]
\[f\left(4\alpha+\frac{4\pi}{3}\right)=2\cos\left(\alpha+\frac\pi2\right)=-2\sin\alpha=-\frac{30}{17},\]
\Needspace{5\baselineskip}
即
\nopagebreak[4]
\[\sin\alpha=\frac{15}{17}.\]
\Needspace{5\baselineskip}
因 $\alpha\in\left[0,\dfrac\pi2\right]$，取正值
\nopagebreak[4]
\[\cos\alpha=\sqrt{1-\frac{225}{289}}=\sqrt{\frac{64}{289}}=\frac8{17}.\]

\Needspace{5\baselineskip}
对第二个条件，
\nopagebreak[4]
\[\frac14\left(4\beta-\frac{2\pi}{3}\right)+\frac\pi6=\beta-\frac\pi6+\frac\pi6=\beta,\]
\Needspace{5\baselineskip}
故
\nopagebreak[4]
\[f\left(4\beta-\frac{2\pi}{3}\right)=2\cos\beta=\frac85,\qquad \cos\beta=\frac45.\]
\Needspace{5\baselineskip}
因 $\beta\in\left[0,\dfrac\pi2\right]$，
\nopagebreak[4]
\[\sin\beta=\sqrt{1-\frac{16}{25}}=\frac35.\]

\Needspace{5\baselineskip}
由两角和的余弦公式，
\nopagebreak[4]
\[\begin{aligned}
\cos(\alpha+\beta)&=\cos\alpha\cos\beta-\sin\alpha\sin\beta\\
&=\frac8{17}\cdot\frac45-\frac{15}{17}\cdot\frac35\\
&=\frac{32}{85}-\frac{45}{85}=-\frac{13}{85}.
\end{aligned}\]

\gkanswer{答案（1）：$A=2$。$f\left(\dfrac\pi3\right)=A\cos\dfrac\pi4=\dfrac{\sqrt2}{2}A=\sqrt2$。}

\gkanswer{答案（2）：$-\dfrac{13}{85}$。由条件得 $\sin\alpha=\dfrac{15}{17}$，$\cos\alpha=\dfrac{8}{17}$，$\cos\beta=\dfrac45$，$\sin\beta=\dfrac35$，再用两角和的余弦公式。}

\end{answer}
\begin{answer}
\gksolhead{17}{}

\textbf{题目}：

某校 $100$ 名学生期中考试语文成绩的频率分布直方图如图所示，其中成绩分组区间是：$[50,60)$，$[60,70)$，$[70,80)$，$[80,90)$，$[90,100]$。
\nopagebreak[4]
\begin{center}% ============================================================
% f2012_guangdong_arts_q17_histogram.tex
% 2012 广东卷（文科）第 17 题 频率分布直方图（TikZ 矢量重绘）
% 分组 [50,60) [60,70) [70,80) [80,90) [90,100]
% 组距 0.04、0.03、0.02、a、a（y 轴为 频率/组距）
% ============================================================
\begin{tikzpicture}[
  x=1cm, y=1cm, line width=0.6pt, font=\footnotesize, >=Stealth
]

  % 坐标轴
  \draw[->] (0,0) -- (7.15,0);
  \node[right] at (7.15,0.12) {成绩};
  \draw[->] (0,0) -- (0,3.75);
  \node[align=center] at (0,4.15) {$\dfrac{\text{频率}}{\text{组距}}$};

  % 纵轴刻度与虚线
  \foreach \y/\lab/\xe in {0.4/a/6.6, 1.6/0.02/5.6, 2.4/0.03/4.6, 3.2/0.04/3.6} {
    \node[left] at (0,\y) {$\lab$};
    \draw[dashed] (0,\y) -- (\xe,\y);
  }

  % 五个矩形（宽 1，自 x=1.6 起）
  \draw (1.6,0) rectangle (2.6,0.4);
  \draw (2.6,0) rectangle (3.6,3.2);
  \draw (3.6,0) rectangle (4.6,2.4);
  \draw (4.6,0) rectangle (5.6,1.6);
  \draw (5.6,0) rectangle (6.6,0.4);

  % 横轴刻度
  \foreach \x/\lab in {1.6/50, 2.6/60, 3.6/70, 4.6/80, 5.6/90, 6.6/100} {
    \node[below] at (\x,0) {$\lab$};
  }
  \node[below] at (0.25,0) {$0$};
  % 断轴记号
  \draw (0.62,-0.10) -- (0.78,0.10);
  \draw (0.92,-0.10) -- (1.08,0.10);

\end{tikzpicture}
\end{center}
\nopagebreak[4]
（1） 求图中 $a$ 的值；

（2） 根据频率分布直方图，估计这 $100$ 名学生语文成绩的平均分；

（3） 若这 $100$ 名学生语文成绩某些分数段的人数（$x$）与数学成绩相应分数段的人数（$y$）之比如下表所示，求数学成绩在 $[50,90)$ 之外的人数。

\begin{center}
\begin{tabular}{lllll}
\toprule
分数段 & $[50,60)$ & $[60,70)$ & $[70,80)$ & $[80,90)$ \\
\midrule
$x:y$ & $1:1$ & $2:1$ & $3:4$ & $4:5$ \\
\bottomrule
\end{tabular}
\end{center}

\Needspace{4\baselineskip}
\textbf{解析}

组距为 $10$，各组频率等于矩形高乘以 $10$。

\textbf{（1）}

\Needspace{5\baselineskip}
五个矩形的高依次为 $a$，$0.04$，$0.03$，$0.02$，$a$，频率之和为 $1$：
\nopagebreak[4]
\[10(a+0.04+0.03+0.02+a)=1,\]
\Needspace{5\baselineskip}
即
\nopagebreak[4]
\[10(2a+0.09)=1,\qquad 2a+0.09=0.1,\]
\Needspace{5\baselineskip}
解得
\nopagebreak[4]
\[a=0.005.\]

\textbf{（2）}

\Needspace{5\baselineskip}
各组频率依次为
\nopagebreak[4]
\[0.05,\quad 0.4,\quad 0.3,\quad 0.2,\quad 0.05.\]
\Needspace{5\baselineskip}
用各组中点 $55$，$65$，$75$，$85$，$95$ 估计平均分：
\nopagebreak[4]
\[\begin{aligned}
\bar x&=55\times0.05+65\times0.4+75\times0.3+85\times0.2+95\times0.05\\
&=2.75+26+22.5+17+4.75\\
&=73.
\end{aligned}\]
估计平均分为 $73$。

\textbf{（3）}

\Needspace{5\baselineskip}
语文各分数段人数为频率乘以 $100$：
\nopagebreak[4]
\[5,\quad 40,\quad 30,\quad 20.\]
按表中 $x:y$ 求数学相应人数：

\begin{center}
\begin{tabular}{llll}
\toprule
分数段 & 语文人数 $x$ & $x:y$ & 数学人数 $y$ \\
\midrule
$[50,60)$ & $5$ & $1:1$ & $5$ \\
$[60,70)$ & $40$ & $2:1$ & $20$ \\
$[70,80)$ & $30$ & $3:4$ & $40$ \\
$[80,90)$ & $20$ & $4:5$ & $25$ \\
\bottomrule
\end{tabular}
\end{center}

\Needspace{5\baselineskip}
这四段合计
\nopagebreak[4]
\[5+20+40+25=90.\]
\Needspace{5\baselineskip}
故数学成绩在 $[50,90)$ 之外的人数为
\nopagebreak[4]
\[100-90=10.\]

\gkanswer{答案（1）：$a=0.005$。各组矩形高之和乘组距 $10$ 等于 $1$，即 $10(2a+0.09)=1$。}

\gkanswer{答案（2）：$73$。各组取中点乘以对应频率再求和。}

\gkanswer{答案（3）：$10$。语文各段人数为 $5,40,30,20$，按比例得数学各段人数 $5,20,40,25$，合计 $90$，故 $[50,90)$ 之外有 $100-90=10$ 人。}

\end{answer}
\begin{answer}
\gksolhead{18}{}

\textbf{题目}：

如图所示，在四棱锥 $P-ABCD$ 中，$AB\perp$ 平面 $PAD$，$AB\parallel CD$，$PD=AD$，$E$ 是 $PB$ 中点，$F$ 是 $DC$ 上的点且 $DF=\dfrac12AB$，$PH$ 为 $\triangle PAD$ 中 $AD$ 边上的高。

（1） 证明：$PH\perp$ 平面 $ABCD$；

（2） 若 $PH=1$，$AD=\sqrt2$，$FC=1$，求三棱锥 $E-BCF$ 的体积；

（3） 证明：$EF\perp$ 平面 $PAB$。
\nopagebreak[4]
\begin{center}% ============================================================
% f2012_guangdong_arts_q18_pyramid.tex
% 2012 广东卷（文科）第 18 题 四棱锥 P-ABCD（TikZ 矢量重绘）
% AB⊥平面 PAD，AB∥CD，PD=AD，E 为 PB 中点，F 在 DC 上且 DF=AB/2，
% PH 为 △PAD 中 AD 边上的高；按原页拓扑示意，不按比例
% ============================================================
\begin{tikzpicture}[
  x=1cm, y=1cm, line width=0.6pt, font=\footnotesize, >=Stealth
]

  \coordinate (A) at (0,0);
  \coordinate (B) at (3.1,0);
  \coordinate (C) at (6.3,1.6);
  \coordinate (D) at (2.7,1.6);
  \coordinate (P) at (1.1,3.7);
  \coordinate (H) at (1.35,0.8);
  \coordinate (E) at (2.1,1.85);
  \coordinate (F) at (4.2,1.6);

  % 可见棱
  \draw (P) -- (A);
  \draw (P) -- (B);
  \draw (P) -- (C);
  \draw (A) -- (B);
  \draw (B) -- (C);

  % 不可见棱与辅助线
  \draw[dashed] (P) -- (D);
  \draw[dashed] (P) -- (H);
  \draw[dashed] (A) -- (D);
  \draw[dashed] (D) -- (C);
  \draw[dashed] (B) -- (F);
  \draw[dashed] (E) -- (F);

  % 标注
  \node[above] at (P) {$P$};
  \node[below left] at (A) {$A$};
  \node[below] at (B) {$B$};
  \node[right] at (C) {$C$};
  \node at (2.33,1.22) {$D$};
  \node at (1.66,2.00) {$E$};
  \node at (4.28,1.95) {$F$};
  \node at (1.22,0.42) {$H$};

\end{tikzpicture}
\end{center}
\nopagebreak[4]
\Needspace{4\baselineskip}
\textbf{解析}

\textbf{（1）}

\Needspace{5\baselineskip}
因为 $AB\perp$ 平面 $PAD$，而 $PH\subset$ 平面 $PAD$，所以
\nopagebreak[4]
\[AB\perp PH.\]
\Needspace{5\baselineskip}
又 $PH$ 是 $\triangle PAD$ 中 $AD$ 边上的高，故
\nopagebreak[4]
\[PH\perp AD.\]
\Needspace{5\baselineskip}
在平面 $ABCD$ 内，$AB$ 与 $AD$ 相交于 $A$，于是
\nopagebreak[4]
\[PH\perp\text{\text{平面} }ABCD.\]

\textbf{（2）}

\Needspace{5\baselineskip}
$E$ 是 $PB$ 的中点，而 $B$ 在平面 $BCF$（即平面 $ABCD$）内，故 $E$ 到平面 $BCF$ 的距离等于 $P$ 到该平面距离的一半：
\nopagebreak[4]
\[V_{E-BCF}=\frac12V_{P-BCF}.\]
\Needspace{5\baselineskip}
由（1），$P$ 到平面 $ABCD$ 的距离就是 $PH=1$，所以
\nopagebreak[4]
\[V_{P-BCF}=\frac13\cdot S_{\triangle BCF}\cdot PH.\]

\Needspace{5\baselineskip}
$F$ 在 $DC$ 上，$FC=1$。因为 $AB\parallel CD$ 且 $AB\perp AD$，所以 $DC\perp AD$，即 $AD$ 是平行线 $AB$ 与 $DC$ 之间的公垂线段。$B$ 在直线 $AB$ 上，故 $B$ 到直线 $DC$（即直线 $FC$）的距离等于 $AD=\sqrt2$。于是
\nopagebreak[4]
\[S_{\triangle BCF}=\frac12\cdot FC\cdot AD=\frac12\times1\times\sqrt2=\frac{\sqrt2}{2}.\]
\Needspace{5\baselineskip}
因此
\nopagebreak[4]
\[\begin{aligned}
V_{E-BCF}&=\frac12\cdot\frac13\cdot\frac{\sqrt2}{2}\cdot1\\
&=\frac{\sqrt2}{12}.
\end{aligned}\]

\textbf{（3）}

取 $AB$ 的中点 $M$，连接 $ME$，$MF$。

\Needspace{5\baselineskip}
因为 $E$，$M$ 分别是 $PB$，$AB$ 的中点，由三角形中位线定理
\nopagebreak[4]
\[ME\parallel PA.\]
\Needspace{5\baselineskip}
又 $AM=\dfrac12AB=DF$，且 $AM\parallel DF$（分别在 $AB$ 与 $DC$ 上），所以四边形 $AMFD$ 是平行四边形，从而
\nopagebreak[4]
\[MF\parallel AD,\qquad MF=AD.\]

\Needspace{5\baselineskip}
由 $AB\perp$ 平面 $PAD$ 得 $AB\perp AD$，$AB\perp PA$，结合 $MF\parallel AD$、$ME\parallel PA$ 得
\nopagebreak[4]
\[AB\perp MF,\qquad AB\perp ME.\]
\Needspace{5\baselineskip}
$MF$ 与 $ME$ 相交于 $M$，故
\nopagebreak[4]
\[AB\perp\text{\text{平面} }MEF,\qquad AB\perp EF.\]

\Needspace{5\baselineskip}
再证 $EF\perp AP$。取 $AP$ 的中点 $K$。在 $\triangle PDA$ 中 $PD=AD$，$K$ 是底边 $PA$ 的中点，故中线也是高：
\nopagebreak[4]
\[DK\perp PA.\]
\Needspace{5\baselineskip}
另一方面，由 $ME\parallel PA$ 且 $ME=\dfrac12PA=AK$ 知四边形 $AKEM$ 是平行四边形，于是 $KE\parallel AM$，且 $KE=AM=\dfrac12AB=DF$；又 $KE\parallel AB\parallel DF$，所以 $KE$ 与 $DF$ 平行且相等，四边形 $KEFD$ 是平行四边形，得
\nopagebreak[4]
\[EF\parallel DK.\]
\Needspace{5\baselineskip}
结合 $DK\perp PA$ 得
\nopagebreak[4]
\[EF\perp PA.\]

\Needspace{5\baselineskip}
于是 $EF\perp AB$，$EF\perp PA$，而 $AB$ 与 $PA$ 相交于 $A$，都在平面 $PAB$ 内，故
\nopagebreak[4]
\[EF\perp\text{\text{平面} }PAB.\]

\gkanswer{答案（1）：见解析版}

\gkanswer{答案（2）：$\dfrac{\sqrt2}{12}$。$E$ 为 $PB$ 中点，故 $V_{E-BCF}=\dfrac12V_{P-BCF}$，而 $S_{\triangle BCF}=\dfrac12\cdot FC\cdot AD=\dfrac{\sqrt2}{2}$，$PH=1$。}

\gkanswer{答案（3）：见解析版}

\end{answer}
\begin{answer}
\gksolhead{19}{}

\textbf{题目}：

设数列 $\{a_n\}$ 的前 $n$ 项和为 $S_n$，数列 $\{S_n\}$ 的前 $n$ 项和为 $T_n$，满足 $T_n=2S_n-n^2$，$n\in\mathbf N^*$。

（1） 求 $a_1$ 的值；

（2） 求数列 $\{a_n\}$ 的通项公式。

\Needspace{4\baselineskip}
\textbf{解析}

\textbf{（1）}

\Needspace{5\baselineskip}
$n=1$ 时，$T_1=S_1=a_1$。代入 $T_n=2S_n-n^2$ 得
\nopagebreak[4]
\[a_1=2a_1-1,\]
\Needspace{5\baselineskip}
解得
\nopagebreak[4]
\[a_1=1.\]

\textbf{（2）}

\Needspace{5\baselineskip}
当 $n\geqslant2$ 时，
\nopagebreak[4]
\[T_n-T_{n-1}=S_n,\]
\Needspace{5\baselineskip}
而由条件
\nopagebreak[4]
\[\begin{aligned}
T_n-T_{n-1}&=\left(2S_n-n^2\right)-\left(2S_{n-1}-(n-1)^2\right)\\
&=2S_n-2S_{n-1}-n^2+(n-1)^2\\
&=2S_n-2S_{n-1}-2n+1.
\end{aligned}\]
\Needspace{5\baselineskip}
与 $T_n-T_{n-1}=S_n$ 比较得
\nopagebreak[4]
\[S_n=2S_n-2S_{n-1}-2n+1,\]
\Needspace{5\baselineskip}
即
\nopagebreak[4]
\[S_n=2S_{n-1}+2n-1.\]
\Needspace{5\baselineskip}
把递推式配成等比形式：设
\nopagebreak[4]
\[S_n+pn+q=2\left(S_{n-1}+p(n-1)+q\right),\]
\Needspace{5\baselineskip}
展开右端为 $2S_{n-1}+(2p-p)n+(-2p+2q-q)$，与 $2S_{n-1}+2n-1$ 比较得
\nopagebreak[4]
\[p=2,\qquad -2p+q=-1,\]
\Needspace{5\baselineskip}
即 $p=2$，$q=3$。于是
\nopagebreak[4]
\[S_n+2n+3=2\left(S_{n-1}+2(n-1)+3\right),\]
\Needspace{5\baselineskip}
数列 $\{S_n+2n+3\}$ 是公比为 $2$ 的等比数列，首项
\nopagebreak[4]
\[S_1+2\times1+3=1+5=6.\]
\Needspace{5\baselineskip}
故
\nopagebreak[4]
\[S_n+2n+3=6\cdot2^{n-1}=3\cdot2^n,\]
\Needspace{5\baselineskip}
即
\nopagebreak[4]
\[S_n=3\cdot2^n-2n-3.\]
（$n=1$ 时右端 $=6-2-3=1=S_1$，对 $n=1$ 也成立。）

\Needspace{5\baselineskip}
当 $n\geqslant2$ 时，
\nopagebreak[4]
\[\begin{aligned}
a_n&=S_n-S_{n-1}\\
&=\left(3\cdot2^n-2n-3\right)-\left(3\cdot2^{n-1}-2(n-1)-3\right)\\
&=3\cdot2^{n-1}-2.
\end{aligned}\]
\Needspace{5\baselineskip}
$n=1$ 时 $3\cdot2^0-2=1=a_1$，故对一切 $n\in\mathbf N^*$
\nopagebreak[4]
\[a_n=3\cdot2^{n-1}-2.\]

\gkanswer{答案（1）：$a_1=1$。$n=1$ 时 $T_1=S_1=a_1$，代入得 $a_1=2a_1-1$。}

\gkanswer{答案（2）：$a_n=3\cdot2^{n-1}-2$。由 $T_n-T_{n-1}=S_n$ 得 $S_n=2S_{n-1}+2n-1$，解得 $S_n=3\cdot2^n-2n-3$，再作差。}

\end{answer}
\begin{answer}
\gksolhead{20}{}

\textbf{题目}：

在平面直角坐标系 $xOy$ 中，已知椭圆 $C_1:\dfrac{x^2}{a^2}+\dfrac{y^2}{b^2}=1$（$a>b>0$）的左焦点 $F_1(-1,0)$，且点 $P(0,1)$ 在 $C_1$ 上。

（1） 求椭圆 $C_1$ 的方程；

（2） 设直线 $l$ 同时与椭圆 $C_1$ 和抛物线 $C_2:y^2=4x$ 相切，求直线 $l$ 的方程。

\Needspace{4\baselineskip}
\textbf{解析}

\textbf{（1）}

\Needspace{5\baselineskip}
点 $P(0,1)$ 在椭圆上，代入得
\nopagebreak[4]
\[\frac{0}{a^2}+\frac1{b^2}=1,\qquad b^2=1.\]
\Needspace{5\baselineskip}
左焦点 $F_1(-1,0)$ 给出 $c=1$。由 $a^2=b^2+c^2$ 得
\nopagebreak[4]
\[a^2=1+1=2.\]
\Needspace{5\baselineskip}
故椭圆方程为
\nopagebreak[4]
\[\frac{x^2}{2}+y^2=1.\]

\textbf{（2）}

\Needspace{5\baselineskip}
先排除斜率不存在的情形：若 $l:x=t$，它与抛物线 $y^2=4x$ 相切只能切于顶点，此时 $t=0$，但 $x=0$ 与椭圆交于 $(0,\pm1)$ 两点，并不相切。故 $l$ 有斜率，设
\nopagebreak[4]
\[l:y=kx+m\quad(k\neq0).\]

\Needspace{5\baselineskip}
与抛物线相切：把 $x=\dfrac{y-m}{k}$ 代入 $y^2=4x$，得
\nopagebreak[4]
\[ky^2-4y+4m=0,\]
\Needspace{5\baselineskip}
相切要求判别式为零：
\nopagebreak[4]
\[\Delta_1=16-16km=0,\qquad km=1.\]

\Needspace{5\baselineskip}
与椭圆相切：代入 $\dfrac{x^2}{2}+y^2=1$ 得
\nopagebreak[4]
\[x^2+2(kx+m)^2=2,\]
\Needspace{5\baselineskip}
即
\nopagebreak[4]
\[(1+2k^2)x^2+4kmx+2m^2-2=0,\]
\Needspace{5\baselineskip}
相切要求
\nopagebreak[4]
\[\Delta_2=(4km)^2-4(1+2k^2)(2m^2-2)=0.\]
\Needspace{5\baselineskip}
化简：
\nopagebreak[4]
\[\begin{aligned}
16k^2m^2&=8(1+2k^2)(m^2-1),\\
2k^2m^2&=(1+2k^2)(m^2-1)=m^2-1+2k^2m^2-2k^2,
\end{aligned}\]
\Needspace{5\baselineskip}
故
\nopagebreak[4]
\[m^2=1+2k^2.\]

\Needspace{5\baselineskip}
把 $m=\dfrac1k$ 代入：
\nopagebreak[4]
\[\frac1{k^2}=1+2k^2,\qquad 1=k^2+2k^4,\]
\Needspace{5\baselineskip}
即
\nopagebreak[4]
\[2k^4+k^2-1=0,\qquad (2k^2-1)(k^2+1)=0.\]
\Needspace{5\baselineskip}
因 $k^2+1>0$，得
\nopagebreak[4]
\[k^2=\frac12,\qquad k=\pm\frac{\sqrt2}{2},\]
相应地 $m=\dfrac1k=\pm\sqrt2$（与 $k$ 同号）。

\Needspace{5\baselineskip}
故
\nopagebreak[4]
\[l:y=\frac{\sqrt2}{2}x+\sqrt2\quad\text{\text{或}}\quad l:y=-\frac{\sqrt2}{2}x-\sqrt2,\]
\Needspace{5\baselineskip}
化为一般式：
\nopagebreak[4]
\[x-\sqrt2y+2=0\quad\text{\text{或}}\quad x+\sqrt2y+2=0.\]

\gkanswer{答案（1）：$\dfrac{x^2}{2}+y^2=1$。$P(0,1)$ 给 $b=1$，$c=1$ 给 $a^2=b^2+c^2=2$。}

\gkanswer{答案（2）：$x-\sqrt2y+2=0$ 或 $x+\sqrt2y+2=0$。设 $l:y=kx+m$，与抛物线相切给 $km=1$，与椭圆相切给 $m^2=1+2k^2$，联立得 $k^2=\dfrac12$。}

\end{answer}
\begin{answer}
\gksolhead{21}{}

\textbf{题目}：

设 $0<a<1$，集合 $A=\{x\in\mathbf R\mid x>0\}$，$B=\{x\in\mathbf R\mid 2x^2-3(1+a)x+6a>0\}$，$D=A\cap B$。

（1） 求集合 $D$（用区间表示）；

（2） 求函数 $f(x)=2x^3-3(1+a)x^2+6ax$ 在 $D$ 内的极值点。

\Needspace{4\baselineskip}
\textbf{解析}

\textbf{（1）}

\Needspace{5\baselineskip}
记 $g(x)=2x^2-3(1+a)x+6a$，其判别式
\nopagebreak[4]
\[\begin{aligned}
\Delta&=9(1+a)^2-48a\\
&=9a^2+18a+9-48a\\
&=9a^2-30a+9=3(3a-1)(a-3).
\end{aligned}\]
因 $0<a<1$，有 $a-3<0$，$\Delta$ 的符号由 $3a-1$ 决定。

\Needspace{5\baselineskip}
情形一：$0<a<\dfrac13$，此时 $\Delta>0$，$g(x)=0$ 有两个不等实根
\nopagebreak[4]
\[x_{1,2}=\frac{3(1+a)\mp\sqrt{9a^2-30a+9}}{4}.\]
\Needspace{5\baselineskip}
两根之和 $\dfrac{3(1+a)}{2}>0$，之积 $\dfrac{6a}{2}=3a>0$，故 $0<x_1<x_2$，且
\nopagebreak[4]
\[B=(-\infty,x_1)\cup(x_2,+\infty).\]
\Needspace{5\baselineskip}
与 $A=(0,+\infty)$ 取交：
\nopagebreak[4]
\[D=(0,x_1)\cup(x_2,+\infty).\]

\Needspace{5\baselineskip}
情形二：$a=\dfrac13$，此时 $\Delta=0$，$g(x)=2\left(x-1\right)^2$，条件 $g(x)>0$ 即 $x\neq1$，故
\nopagebreak[4]
\[B=\mathbf R\setminus\{1\},\qquad D=(0,1)\cup(1,+\infty).\]

\Needspace{5\baselineskip}
情形三：$\dfrac13<a<1$，此时 $\Delta<0$，$g(x)>0$ 对一切实数成立，$B=\mathbf R$，故
\nopagebreak[4]
\[D=(0,+\infty).\]

\Needspace{5\baselineskip}
综上，
\nopagebreak[4]
\[D=\begin{cases}
(0,x_1)\cup(x_2,+\infty), & 0<a<\dfrac13,\\[1mm]
(0,1)\cup(1,+\infty), & a=\dfrac13,\\[1mm]
(0,+\infty), & \dfrac13<a<1,
\end{cases}\]
其中 $x_{1,2}=\dfrac{3(1+a)\mp\sqrt{9a^2-30a+9}}4$。

\textbf{（2）}

\Needspace{5\baselineskip}
求导：
\nopagebreak[4]
\[\begin{aligned}
f'(x)&=6x^2-6(1+a)x+6a\\
&=6\left(x^2-(1+a)x+a\right)=6(x-1)(x-a).
\end{aligned}\]
因 $0<a<1$，有 $a<1$，故 $f'(x)$ 的符号为：

\begin{itemize}
\item $x<a$ 时 $f'(x)>0$，$f$ 单调递增；
\item $a<x<1$ 时 $f'(x)<0$，$f$ 单调递减；
\item $x>1$ 时 $f'(x)>0$，$f$ 单调递增。
\end{itemize}

所以 $x=a$ 是极大值点，$x=1$ 是极小值点，但极值点必须落在 $D$ 内。

由 $g(a)=2a^2-3(1+a)a+6a=-a^2+3a=a(3-a)>0$（$0<a<1$）知 $a\notin[x_1,x_2]$；又 $g(1)=3a-1<0$ 说明 $1\in(x_1,x_2)$，而 $a<1$，故 $a<x_1$，即 $a\in(0,x_1)\subset D$。

逐个情形判断：

\begin{itemize}
\item $0<a<\dfrac13$：$D=(0,x_1)\cup(x_2,+\infty)$。此时 $a\in(0,x_1)\subset D$，而 $g(1)=3a-1<0$ 说明 $1\in(x_1,x_2)$，即 $1\notin D$。故 $D$ 内只有极大值点 $x=a$。
\item $a=\dfrac13$：$D=(0,1)\cup(1,+\infty)$。此时 $a=\dfrac13\in(0,1)\subset D$，而 $1\notin D$。故 $D$ 内只有极大值点 $x=\dfrac13$。
\item $\dfrac13<a<1$：$D=(0,+\infty)$，$a$ 与 $1$ 都在 $D$ 内。故 $D$ 内极大值点为 $x=a$，极小值点为 $x=1$。
\end{itemize}

（补充：$x=a$ 处极大值为 $f(a)=-a^3+3a^2$，$x=1$ 处极小值为 $f(1)=3a-1$。）

\gkanswer{答案（1）：$0<a<\dfrac13$ 时 $D=(0,x_1)\cup(x_2,+\infty)$，其中 $x_{1,2}=\dfrac{3(1+a)\mp\sqrt{9a^2-30a+9}}4$；$a=\dfrac13$ 时 $D=(0,1)\cup(1,+\infty)$；$\dfrac13<a<1$ 时 $D=(0,+\infty)$。按判别式 $\Delta=3(3a-1)(a-3)$ 的符号分三种情形。}

\gkanswer{答案（2）：$0<a<\dfrac13$ 时只有极大值点 $x=a$；$a=\dfrac13$ 时只有极大值点 $x=\dfrac13$；$\dfrac13<a<1$ 时极大值点 $x=a$、极小值点 $x=1$。由 $f'(x)=6(x-1)(x-a)$ 及 $a$、$1$ 是否落在 $D$ 内判定。}

\end{answer}